\documentclass[11pt,a4paper]{article}
\usepackage[margin=2.4cm]{geometry}
\usepackage{amsmath,amssymb,booktabs,hyperref}
\newcommand{\avg}[1]{\{\!\{#1\}\!\}}
\newcommand{\dt}{\Delta t}
\newcommand{\code}[1]{\texttt{#1}}
\title{The fast system of the 2D DGSEM:\\Verlet-type small steps and the column solve by static condensation}
\author{MultiRate notes}
\date{\today}
\begin{document}
\maketitle

\begin{abstract}
The linear fast system of the 2D DGSEM is written in two groups of unknowns, the thermodynamic variables $T$ and the momentum
$M$, with the coupling operators $A_T$, $A_M$ and the self-couplings $D_T$, $D_M$, each split into the horizontal ($x$, $\xi$) and the vertical
($z$, $\eta$) part, for the DGSEM with LMARS in $\rho\mathbf u$. All on a Cartesian grid, without metric terms. The fast options are
written as schemes; the Verlet-type schemes have three stages. The column system of the vertically implicit scheme is solved as in
hybridized DG: the interface momentum is introduced only in the linear algebra, the element unknowns are eliminated locally, and a
tridiagonal system in the interface momenta remains per node column.
\end{abstract}

\section{Setting}
\paragraph{Fast system.} Coefficients frozen at the stage: $\theta$, $A=\mathrm dp/\mathrm d\Theta=\gamma p/\Theta$, $c^2=A\theta$;
\[
\rho_t+\nabla\cdot m=0,\qquad \Theta_t+\nabla\cdot(\theta m)=0,\qquad p_t=A\,\Theta_t,\qquad m_t+\nabla p+g\rho\,\mathbf e_z=0 .
\]
Groups: thermodynamic variables $T=(\rho,\Theta,p)$, momentum $M=(m_x,m_z)$.
The slow part enters as a forcing $f=(f_T,f_M)$, constant during a MIS stage.

\paragraph{Grid.} Cartesian elements of size $\Delta x_e\times\Delta z_e$, tensor LGL nodes of degree $P$, collocated quadrature, SBP derivatives $D_x$, $D_z$
along the node lines (scaled by $2/\Delta x_e$, $2/\Delta z_e$). $x$ ($\xi$) is horizontal, $z$ ($\eta$) vertical. Faces: $x$-faces ($x=$const, between horizontal
neighbours) and $z$-faces ($z=$const, between vertical neighbours); walls at the ground and the top are $z$-faces. At a face node of
element $K$ with outward unit normal $\nu$ ($\pm\mathbf e_x$ on $x$-faces, $\pm\mathbf e_z$ on $z$-faces): $m_\nu=\nu\cdot m$, the inner value $(\cdot)_{\rm in}$ of $K$ and the
outer value $(\cdot)_{\rm out}$ of the neighbour, $\avg x=\tfrac12(x_{\rm in}+x_{\rm out})$, and the lifting factor
\[
\sigma^x=\frac{2}{\Delta x_e\,w_1},\qquad \sigma^z=\frac{2}{\Delta z_e\,w_1}\qquad(w_1=2/(P(P+1))\text{ the end weight}).
\]
$c$ is the frozen face sound speed, $\theta_{\rm up}$ the upwind value.

\paragraph{Directions.} On the Cartesian grid every $x$-term acts only on $m_x$ and every $z$-term only on $m_z$: the horizontal operators
couple $T$ with $m_x$, the vertical ones $T$ with $m_z$. A node column (the nodes of one $x$-position in a column of elements) couples
through the $z$-terms only.

\paragraph{Operators as matrices.} $T$ and $M$ denote vectors of all node values, ordered element by element. An operator such as $A_T^z$ is the matrix that maps the momentum vector to the
$z$-part of the tendency of $T$; its row index runs over the nodes and variables of $T$, its column index over those of $M$. A block of
such a matrix is labelled by the pair (receiving unknown, acting unknown): an \emph{element block} couples the nodes of one element (volume
terms, and the face terms that use only the element's own values), a \emph{neighbour block} couples a face node of one element with the
face node of the neighbour across the face (face terms with $\avg\cdot$ and jumps). $D_T$, $D_M$ are the parts of the tendency of a group that act on the same group.

\paragraph{MIS stage and small step.} Stage $i$ of a MIS step integrates the linear fast system with the forcing $f^{(i)}$ over $[0,d_i\dt]$ with
$n_i$ small steps of length $h=d_i\dt/n_i$. A small step maps $(T^0,M^0)$ to $(T^1,M^1)$; $T^{1/2}$ is the value after the first half step. Parameter: off-centring $\beta=(1+\varepsilon)/2$ of the
vertical force (no extrapolation, no divergence damping).

\section{DGSEM with LMARS in $\rho\mathbf u$}
\paragraph{LMARS.} Seen from $K$, with the characteristic interface pressure $P^*=p_{\rm in}-c\,(m^*_\nu-m_{\nu,\rm in})=p_{\rm out}+c\,(m^*_\nu-m_{\nu,\rm out})$ of the two
sides:
\[
m^*_\nu=\avg{m_\nu}-\frac{p_{\rm out}-p_{\rm in}}{2c},\qquad p^*=\avg p-\frac c2\,(m_{\nu,\rm out}-m_{\nu,\rm in});\qquad\text{walls: }m^*_\nu=0,\ p^*=p_{\rm in}+c\,m_{\nu,\rm in}.
\]

\paragraph{Fast system.}
\begin{equation}
T_t=A_T\,M+D_T\,T+f_T,\qquad M_t=A_M\,T+D_M\,M+f_M,\label{eq:std}
\end{equation}
$A_T=A_T^x+A_T^z$, $A_M=A_M^x+A_M^z$, $D_T=D_T^x+D_T^z$, $D_M=D_M^x+D_M^z$. With $d\in\{x,z\}$ and the face terms at the face nodes of the
$d$-faces:
\begin{itemize}
\item $A_T^d$ ($m_d\to T$), the central part of the mass flux:
  \[
  \rho_t\ni-D_d\,m_d-\sigma^d\big(\avg{m_\nu}-m_{\nu,\rm in}\big),\qquad \Theta_t\ni-D_d(\theta m_d)-\sigma^d\big(\theta_{\rm up}\avg{m_\nu}-\theta\,m_{\nu,\rm in}\big),
  \]
  $p_t\ni A\times(\Theta_t)$; at walls $\rho_t\ni\sigma^z m_{\nu,\rm in}$, $\Theta_t\ni\sigma^z\theta\,m_{\nu,\rm in}$.
\item $A_M^d$ ($T\to m_d$), pressure gradient, central face pressure, gravity in $z$:
  \[
  (m_d)_t\ni-D_d\,p-\sigma^d\,\nu_d\big(\avg p-p_{\rm in}\big),\qquad (m_z)_t\ni-g\rho\quad(d=z).
  \]
\item $D_T^d$ ($T\to T$), the pressure jump in $m^*$:
  \[
  \rho_t\ni\sigma^d\,\frac{p_{\rm out}-p_{\rm in}}{2c},\qquad \Theta_t\ni\sigma^d\,\theta_{\rm up}\frac{p_{\rm out}-p_{\rm in}}{2c},\qquad p_t\ni A\times(\Theta_t).
  \]
\item $D_M^d$ ($m_d\to m_d$), the momentum jump in $p^*$ and the wall term:
  \[
  (m_d)_t\ni\sigma^d\,\nu_d\,\frac c2\,(m_{\nu,\rm out}-m_{\nu,\rm in});\qquad\text{walls: }(m_z)_t\ni-\sigma^z\,\nu_z\,c\,m_{\nu,\rm in}.
  \]
\end{itemize}
$D_T$ and $D_M$ couple only the two nodes of a face node pair. Both groups have a self-coupling.

\paragraph{Schemes.}

\medskip\noindent\textbf{S1 (explicit Störmer--Verlet).}
\begin{enumerate}
\item $T^{1/2}=T^0+\tfrac h2\big(A_TM^0+D_TT^0+f_T\big)$
\item $M^1=M^0+h\big(A_MT^{1/2}+D_MM^0+f_M\big)$
\item $T^1=T^{1/2}+\tfrac h2\big(A_TM^1+D_TT^{1/2}+f_T\big)$
\end{enumerate}
Explicit in $x$ and $z$; the step is limited by the vertical node spacing.

\medskip\noindent\textbf{S2 (Störmer--Verlet, $z$ implicit; Verwer's component splitting).}
\begin{enumerate}
\item $T^{1/2}=T^0+\tfrac h2\big(A_T^xM^0+D_T^xT^0+f_T\big)+\tfrac h2\big(A_T^zM^0+D_T^zT^0\big)$
\item $M^1=M^0+h\big(A_M^xT^{1/2}+D_M^xM^0+f_M\big)+h\big[(1-\beta)\big(A_M^zT^0+D_M^zM^0\big)+\beta\big(A_M^zT^1+D_M^zM^1\big)\big]$
\item $T^1=T^{1/2}+\tfrac h2\big(A_T^xM^1+D_T^xT^{1/2}+f_T\big)+\tfrac h2\big(A_T^zM^1+D_T^zT^1\big)$
\end{enumerate}
Stage 1 is explicit in both directions: its horizontal part ($A_T^xM^0$, $D_T^xT^0$) and its vertical part ($A_T^zM^0$, $D_T^zT^0$) use the
old values; together with the implicit vertical part of stage 3 the vertical part of $T$ is integrated by the trapezoidal rule. Stages 2 and 3 are coupled through the $z$-terms
($T^1$, $m_z^1$); the linear equation that remains is derived below.

\paragraph{From stages 2 and 3 to the column system (S2).} Without metric, $A_M^x$ and $D_M^x$ act only on $m_x$, and $A_M^z$, $D_M^z$ only on $m_z$. Stage 2
therefore splits into its two components:
\begin{align*}
m_x^1&=m_x^0+h\big(A_M^xT^{1/2}+D_M^xm_x^0+f_{m_x}\big)\qquad\text{(explicit)},\\
m_z^1&=m_z^*+\beta h\big(A_M^zT^1+D_M^zm_z^1\big),\qquad m_z^*=m_z^0+h f_{m_z}+(1-\beta)h\big(A_M^zT^0+D_M^zm_z^0\big).
\end{align*}
With $m_x^1$ known, the horizontal part of stage 3 is explicit as well, and stage 3 reads
\[
T^1=T^*+\tfrac h2\big(A_T^zm_z^1+D_T^zT^1\big),\qquad T^*=T^{1/2}+\tfrac h2\big(A_T^xm_x^1+D_T^xT^{1/2}+f_T\big).
\]
The two implicit lines form one linear system,
\[
\begin{pmatrix}I-\tfrac h2D_T^z&-\tfrac h2A_T^z\\-\beta h\,A_M^z&I-\beta h\,D_M^z\end{pmatrix}\begin{pmatrix}T^1\\m_z^1\end{pmatrix}=\begin{pmatrix}T^*\\m_z^*\end{pmatrix}.
\]
All its blocks are $z$-operators: it decouples into one system per node column. Eliminating $T^1=B^{-1}\big(T^*+\tfrac h2A_T^zm_z^1\big)$, $B=I-\tfrac h2D_T^z$
($D_T^z$ couples only the two nodes of a $z$-face, so $B$ is local), gives the equation for the vertical momentum,
\[
\Big(I-\beta h\,D_M^z-\tfrac{\beta h^2}{2}\,A_M^zB^{-1}A_T^z\Big)\,m_z^1=m_z^*+\beta h\,A_M^zB^{-1}T^*,
\]
and then $T^1=B^{-1}\big(T^*+\tfrac h2A_T^zm_z^1\big)$. $A_M^zB^{-1}A_T^z$ is the vertical pressure force of the vertical mass flux (a second
difference along the column, gravity included); $D_M^z$ the momentum damping of the $z$-faces and walls. The system matrix is fixed during a
MIS stage. In the standard formulation $A_T^z$, $A_M^z$ and $D_M^z$ contain neighbour blocks (the central face values and the jumps couple the
two elements at a $z$-face), so the column matrix is block tridiagonal over the elements of the column, with blocks of the size of the
element unknowns; in this form there is no local elimination (Section~\ref{sec:hdg} rewrites it with the interface momentum).
The third stage of a small step and the first of the next use the same $M^1$; combined, only the first and the last half step of a MIS stage
remain, but the horizontal self-coupling has to be evaluated with the new value, $D_T^xT^1$.

\medskip\noindent\textbf{S3 (forward--backward, $z$ implicit).}
\begin{enumerate}
\item $M^1=M^0+h\big(A_M^xT^0+D_M^xM^0+f_M\big)+h\big[(1-\beta)\big(A_M^zT^0+D_M^zM^0\big)+\beta\big(A_M^zT^1+D_M^zM^1\big)\big]$
\item $T^1=T^0+h\big(A_T^xM^1+D_T^xT^0+f_T\big)+h\big(A_T^zM^1+D_T^zT^1\big)$
\end{enumerate}
The column system follows as for S2 with the weight $h$ in place of $\tfrac h2$ on $T$: $\big(I-\beta hD_M^z-\beta h^2A_M^zB^{-1}A_T^z\big)m_z^1=\dots$, $B=I-hD_T^z$;
$T$ and $M$ on the same level.

\section{The column system of S2 by static condensation}\label{sec:hdg}
The interface momentum is used here only to solve the column system of S2; the scheme and its solution are those of the
standard formulation.

\paragraph{The column system.} One node column with the elements $K=1,\dots,N$ from bottom to top and the interior $z$-faces
$f=1,\dots,N-1$ between $K=f$ and $K=f+1$ (one node pair per face in the column). The system of S2 is
\[
\begin{pmatrix}I-\tfrac h2D_T^z&-\tfrac h2A_T^z\\-\beta h\,A_M^z&I-\beta h\,D_M^z\end{pmatrix}\begin{pmatrix}T^1\\m_z^1\end{pmatrix}=\begin{pmatrix}T^*\\m_z^*\end{pmatrix},
\]
with neighbour blocks at every $z$-face (the central values $\avg{m_\nu}$, $\avg p$ and the jumps).

\paragraph{Interface momentum.} On every interior $z$-face node an auxiliary unknown $\hat M_f$ is introduced with the exact constraint
\[
\hat M_f=m^*=\avg{m_n}-\frac{p_R-p_L}{2c}\qquad(n\text{ from }L=K_f\text{ to }R=K_{f+1}),
\]
equivalent to $P^*_L=P^*_R$. Since $p^*=P^*=p_{\rm in}-c\,(m^*_\nu-m_{\nu,\rm in})$, all face terms of an element are written with its own values
and $\hat M$:
\[
\rho_t\ni-\sigma^z(\hat M_\nu-m_{\nu,\rm in}),\qquad \Theta_t\ni-\sigma^z(\theta_{\rm up}\hat M_\nu-\theta\,m_{\nu,\rm in}),\qquad p_t\ni A\times(\Theta_t),\qquad
(m_z)_t\ni\sigma^zc\,\nu_z(\hat M_\nu-m_{\nu,\rm in}).
\]
With $\hat M=m^*$ these are the face terms of $A_T^z$, $D_T^z$, $A_M^z$, $D_M^z$; the operator is unchanged.

\paragraph{Blocks.} Let $Z_K=(T_K,m_K)$ be the unknowns of element $K$ in the column ($\rho,\Theta,p,m_z$ at its $P+1$ nodes). The column system
becomes
\[
E_K\,Z_K-B_K\,\hat M=R_K\quad(K=1,\dots,N),\qquad \hat M_f=\sum_KG_{fK}Z_K\quad(f=1,\dots,N-1),
\]
\begin{itemize}
\item $E_K=I-\Delta\,J_{KK}$, $\Delta=\mathrm{diag}(\tfrac h2\ \text{on }T,\ \beta h\ \text{on }m_z)$, $J_{KK}$ the vertical operator restricted to element $K$:
  the volume terms ($-D_zm_z$, $-D_z(\theta m_z)$, $-D_zp-g\rho$) and the face terms with the element's own values ($+\sigma^zm_{\nu,\rm in}$,
  $+\sigma^z\theta\,m_{\nu,\rm in}$, $-\sigma^zc\,m_{\nu,\rm in}$; walls as before);
\item $B_K=\Delta\,J_{K\hat M}$, the face terms with $\hat M$; element $K$ sees $\hat M_{K-1}$ (bottom) and $\hat M_K$ (top);
\item $G_{fK}$, the constraint: $\tfrac12 m_n$ and $\pm p/(2c)$ at the top node of $K=f$ and at the bottom node of $K=f+1$;
\item $R_K=(T^*_K,m^*_K)$.
\end{itemize}
Different elements are coupled only through $\hat M$: there are no neighbour blocks.

\paragraph{Reduction.}
\begin{enumerate}
\item \emph{Elements}: $Z_K=y_K+W_K\hat M$, $y_K=E_K^{-1}R_K$, $W_K=E_K^{-1}B_K$; one dense system of size $4(P+1)$ per element, $W_K$ with the
  two columns of its faces.
\item \emph{Interface momenta}: the constraint with $Z_K$ inserted,
  \[
  \Big(I-\sum_KG_{\cdot K}W_K\Big)\hat M=\sum_KG_{\cdot K}\,y_K .
  \]
  Face $f$ couples through the elements $f$ and $f+1$ with the faces $f-1$, $f$, $f+1$: the matrix is tridiagonal, of size $N-1$ per node
  column.
\item \emph{Back substitution}: $Z_K=y_K+W_K\hat M$.
\end{enumerate}
$\hat M$ is not kept after the solve. The result is the solution of the column system of S2 (static condensation of hybridized DG:
the element unknowns are eliminated locally and only the interface unknowns are coupled). On terrain the same holds with $m_x$ in $Z_K$
and the contravariant $\hat M=\tilde n\cdot m^*$.

\paragraph{Implementation.} \code{DGPressureModel(base; lmars = :momentum, vsolver = :hdg)} (\code{src/DGSEM/Euler2D/frozen\_p.jl},
\code{\_hdg\_factor}): $J_{KK}$, $J_{K\hat M}$ and $G$ are taken from the vertical operator written with $\hat M$; \code{vsolver = :lu} solves the same
system by a sparse LU.

\section{Remarks}
\begin{itemize}
\item S1, S2 are symmetric for $\varepsilon=0$ and $D^x=0$; an explicit horizontal self-coupling ($D_T^x$, $D_M^x$) makes them of first order.
  S3 is of first order.
\item The $z$-parts are unconditionally stable for $\beta\ge\tfrac12$; the $x$-parts are limited by the Störmer--Verlet limit of the horizontal operator,
  reduced by the explicit $D^x$.
\item A state with $T_t=M_t=0$ is a fixed point of every stage of every scheme: steady states are kept.
\item On terrain the $x$- and $z$-operators also couple $m_x$ and $m_z$ through the metric; that is the next step.
\end{itemize}
\end{document}
